% !TEX program = xelatex
\documentclass[12pt]{article}
\usepackage[margin=28mm]{geometry}
\usepackage{hyperref}
\usepackage{amsmath, amsthm, amssymb, amsfonts}
\usepackage{tikz-cd}
\usepackage{fontspec}
\usepackage{unicode-math}
\setmainfont{Latin Modern Roman}
\setmonofont{Latin Modern Mono}
\usepackage{listings}
\lstset{basicstyle=\ttfamily\footnotesize, columns=fullflexible, breaklines=true, showstringspaces=false}

\title{lean\_advanced\_task — Lean→TeX Export (English)}
\author{TeX generated by ChatGPT-5 Pro\\Lean source generated by CODEX (OpenAI)}
\date{2025-09-10}

\begin{document}
\maketitle

\begin{abstract}
This document was automatically generated from a Lean source file. It includes provenance, sample diagrams, a quick stats table, a declarations table (first-line signatures), and the full original Lean listing.
\end{abstract}

\section*{Provenance}
This \LaTeX{} file was generated from the Lean source file \texttt{lean\_advanced\_task.lean}.\\
\textbf{TeX generation}: ChatGPT-5 Pro.\\
\textbf{Lean source generation}: CODEX (OpenAI).

\section*{At-a-glance}
Total lines: 22\\
Defs: 0\\
Lemmas: 0\\
Theorems: 1\\
Structures: 0\\
Inductives: 0

\section*{Sample diagrams}

\[
\begin{tikzcd}
A \arrow[r, "f"] \arrow[d, "g"'] & B \arrow[d, "h"] \\
C \arrow[r, "k"'] & D
\end{tikzcd}
\]



\[
\begin{tikzcd}
T \arrow[r, "h"] \arrow[dr, swap, "\pi_1\circ h"] \arrow[drr, "\pi_2\circ h"] & X\times Y \arrow[d, swap, "\pi_1"] \arrow[dr, "\pi_2"] \\
& X & Y
\end{tikzcd}
\]


\section*{Declarations (first-line signatures)} 
\begin{center}
\begin{tabular}{lll}
\textbf{kind} & \textbf{name} & \textbf{signature (first line)} \\ \hline
theorem & \texttt{linear\_map\_add\_property} & \texttt{{R V W : Type*} [Semiring R] [AddCommMonoid V] [AddCommMonoid W] [Module R V] [Module R W] (f : V →ₗ[R] W) (v w : V) : f (v + w) = f v + f w := by} \\
\end{tabular}
\end{center}

\section*{Lean source (verbatim)}
\begin{lstlisting}
import Mathlib.Algebra.Module.LinearMap.Defs
import Mathlib.Algebra.Module.LinearMap.Basic
import Mathlib.Algebra.Module.Prod

-- 線形写像の加法性の証明
theorem linear_map_add_property {R V W : Type*} 
  [Semiring R] [AddCommMonoid V] [AddCommMonoid W] 
  [Module R V] [Module R W]
  (f : V →ₗ[R] W) (v w : V) : 
  f (v + w) = f v + f w := by
  simpa using (LinearMap.map_add f v w)

-- Bourbaki スタイルの軽い example：順序対と射影で成分を読み取る
section Examples
  variable {R V W : Type*}
  variable [Semiring R] [AddCommMonoid V] [AddCommMonoid W]
  variable [Module R V] [Module R W]
  example (f : V →ₗ[R] W) (p : V × V) : f (p.1 + p.2) = f p.1 + f p.2 := by
    simpa using (LinearMap.map_add f p.1 p.2)
  example (f : V →ₗ[R] W) (v w : V) : f (v + w) = f v + f w := by
    simpa using (LinearMap.map_add f v w)
end Examples

\end{lstlisting}

\end{document}
